% commentaire.tex — f164r, vv. 123–164 (conversion initiale des travaux antérieurs)
\comm{130}{ὦ πότνιαι Γενετυλλίδες.}{Déesses de la génération, les Génétyllides sont aussi citées par Aristophane au vers 52 des \italique{Nuées} et au second vers de \italique{Lysistrata} ; de ce dernier passage on peut inférer que les femmes étaient réputées célébrer des rites licencieux en leur honneur.}
\comm{131}{θηλυδριῶδες.}{Semble être un néologisme d'Aristophane. Chantraine\nc{\italique{Dictionnaire étymologique de la langue grecque}, \italique{sub voce} θηλή.} suppose un relais *θηλύδριον qui expliquerait aussi θηλυδρίας, bien attesté. Aristophane a-t-il porté son choix sur θηλυδριώδης pour faire entendre θηριώδης et toute l'animalité du spasme de Mnésiloque ? (cf. κατεγλωττισμένον et μανδαλωτόν.)}
\comm{131–2}{καὶ κατεγλωττισμένον / καὶ μανδαλωτόν}{: ces mots renvoient tous deux à ce que les anglophones appellent – ne boudons pas l'hommage qui nous est ainsi rendu ! – \italique{french kiss}. À propos de μανδαλωτόν, je ne résiste pas à citer Taillardat (§195) qui met au jour la bigarrure de la langue aristophanienne : « le substantif, τὸ φίλημα, \italique{le baiser}, est déjà connu d'Eschyle ; il se trouve chez Sophocle, chez Euripide, chez les comiques Platon et Nicophon. En revanche, il n'est pas attesté dans les œuvres conservées d'Aristophane, bien que le poëte emploie φιλεῖν comme synonyme de κυνεῖν [...] Aristophane préfère des mots imagés et comme ``techniques'' : τὸ περιπεταστόν (\italique{Acharniens} 1201), τὸ ἐπιμανδαλωτόν, \italique{ib.}, τὸ μανδαλωτόν (\italique{Thesm.} 132), les deux derniers termes évoquant l'idée d'une porte fermée ; ὁ μάνδαλος, comme βάλανος, désigne la clavette assurant la barre de la porte. La clavette, c'est ici la langue. »}
\comm{134}{νεανίσχ᾽ εἴ τις εἶ.}{Le manuscrit de Ravenne comporte une coquille : le κ de νεανίσκ᾽ est aspiré alors qu'il est suivi de εἴ. Plusieurs solutions ont été apportées : lire ἥτις (justifiant le féminin par Hésychius (η 948) glosant ἥτις· ὅστις). D'autres se sont contentés de substituer un κ au χ. J'ai une préférence pour cette dernière correction\nc{Signalons qu'une coquille analogue (οὐχ᾽ εἴρηχ᾽ R en place d'un οὐκ εἴρηχ᾽) se retrouve au vers 556. De plus, les anciens manuscrits sur lesquels on a copié les \italique{Thesophories} n'étaient pas systématiquement élidés ; ce genre de fautes s'explique donc aisément. Cf. par exemple v. 148 ἐσθητ' ἅμα ; v. 279 θραττ' έπου ; v. 1178 ἔρχετ' ὡς.} qui me semble être une pique contre Agathon : « jeune homme... si tu en es un ! » car Agathon serait âgé au moment de la représentation de 35 ans ou environ ; à croire que son mode de vie le conserve bien... mais tout-de-même ! cette allusion permanente à la jeunesse d'Agathon – comme dans le \italique{Banquet} de Platon censé se passer en 416, cinq ans plus tôt – me paraît plutôt être une allusion au fait qu'Agathon pèche par certain attribut pour être tout-à-fait considéré comme un homme\nc{καὶ ποῦ πέος; demande Mnésiloque au vers 142.}.}
\comm{134–5}{κατ᾽ Αἰσχύλον / ἐκ τῆς Λυκουργείας.}{D'après le scholiaste, la \italique{Lycurgie} est le nom collectif d'une tétralogie d'Eschyle, composée des \italique{Édoniens} (fr.\nc{Les fragments d'Eschyle sont cités selon l'édition de S. Radt, \italique{Tragicorum Graecorum Fragmenta}, vol. 3 (1985).} 57-67), des \italique{Bassarides} (fr. 23-25), des \italique{Jouvenceaux} (Νεανίσκοι, fr. 146-9) et du drame satyrique \italique{Lycurgue} (fr. 124-6). Elle compte les mésaventures du roi Lycurgue essayant de contrer l'invasion du culte dionysiaque en Thrace. Laissons la parole à François Jouan\nc{On lira avec intérêt deux de ses articles : « Dionysos chez Eschyle » (in \italique{Kernos} 5, 1992, p. 71-86) et « Une introduction au mythe des \italique{Bacchantes} » (\italique{Kentron} 14, 1998). C'est le premier que nous citons.} : « La légende de Lycurgue était déjà connue d'Homère. L'\italique{Iliade} rappelait sommairement que Lycurgue avait été aveuglé par Zeus pour avoir pourchassé les Nourrices de Dionysos et forcé le petit dieu à se jeter à la mer\nc{\italique{Iliade}, VI, 130-140.}. La version eschyléenne est très différente, et comporte en particulier une localisation précise : le héros est roi des Édoniens et vit en Thrace, dans la région du Pangée. On admet en général que cette tétralogie [...] datait de la fin de la vie d'Eschyle, entre 466 et 459. Or dans cette période, des forces athéniennes commandées par Cimon guerroyaient précisément contre les Édoniens, autour d'Eion et d'Ennéahodoi. [...] Les \italique{Édoniens} sont la pièce la moins mal connue, en particulier du fait de l'existence de fragments d'un \italique{Lycurgue} du poète latin Naevius, qui semble avoir suivi d'assez près le drame d'Eschyle. Dans le prologue, un messager annonçait au roi thrace l'arrivée d'un personnage entouré d'un bruyant cortège de musiciens et de danseurs – Dionysos et son thiase. Le roi ordonnait de les arrêter et de les faire comparaître devant lui. Dans la parodos, le chœur, sans doute composé de sujets de Lycurgue, décrivait le cortège dionysiaque. Un garde rendait compte de l'arrestation des trublions et de leur docilité. Le roi interrogeait lui-même Dionysos, comme le Penthée des \italique{Bacchantes}. Aussi mécontent des réponses de Dionysos que le prince thébain, il l'emprisonnait avec sa suite dans son palais, où il devait trouver, lui disait-il ironiquement «~un logement gratuit [ἄδασμος, fr. 63]~». Mais la défaite du dieu n'était que provisoire. Bientôt les liens des captifs tombaient d'eux-mêmes, Dionysos apparaissait dans sa gloire sur le toit du palais, qui vacillait sur sa base et se nimbait de flammes : autre motif inclus dans les \italique{Bacchantes}\nc{Cf. Hérodote, II, 145.}. » \\
C'est donc bien à cette première tragédie des \italique{Édoniens} que semble se référer Mnésiloque, interrogeant à la manière de Lycurgue, associant ainsi Agathon à Dionysos (métaphore qui avait commencée d'être filée au vers 41 quand le serviteur d'Agathon parlait du thiase de Muses réunies chez son maître). Rappelons que, pour les Grecs de l'époque, Dionysos était un dieu jeune : « Les particularités de son physique et de son costume sont détaillées aussi bien par Eschyle\nc{Voir les fragments 59, 60, 62 des \italique{Édoniens}.} que par Euripide. Il est grand et mince avec de longues jambes, une chevelure blonde, abondante et parfumée, de belles boucles. Il est beau, avec un teint vermeil et des traits de visage plutôt féminins, l'air souriant. Il est vêtu d'une robe asiatique de fin tissu, tombant jusqu'à ses pieds et tient en main le thyrse. »\nc{F. Jouan, \italique{Une introduction au mythe des} Bacchantes, p. 12.} 
(Ajoutons à propos de cette \italique{crocote} au cœur des péripéties de Mnésiloque qu'elle est aussi associée à Dionysos dans les \italique{Grenouilles} (v.46) ; par Cratinos, fr. 40 de l'édition de Kassel-Austin et, comme on l'a vu, Eschyle, fr. 59.) 
Voilà qui explique un peu les associations d'idées de Mnésiloque ! \nc{Eur. \italique{Bacch}., 586-599 ; 622-624 ; 633.}}
\comm{139–44}{Fragments d'un ancien papyrus.}{L'on possède quelques fragments d'un ancien papyrus (\PsiUn \ dans notre apparat) du premier ou second siècle ap. J.-C. pour les vers 133-144 ; 145-156 ; 237-245 ; 273-288 ; 594-596 ; 804-809. Il nous apprend en outre le \latin{numerus versuum} d'une ancienne édition des \italique{Thesmophories} : la pièce comportait alors 1201 vers (contre 1231 pour celle de Brunck sur laquelle se basent tous les éditeurs modernes). La papyrus est reproduit en annexe. Cf. \italique{Papiri della Società Italiana} XI, n° 1194 et  \italique{Papiri della Società Italiana} XIV, p. XV.}
\comm{139}{τί λήκυθος καὶ στρόφιον;}{Dans sa liste des attributs masculins et féminins qui jurent ensemble chez Agathon, Mnésiloque oppose ici le \italique{lécythe} et le \italique{strophion}. Si le cas du strophion ne pose pas de problème : il s'agit du \latin{strophium} ou \latin{fascia pectoralis}\nc{Je renvoie ici aux articles correspondants de la monumentale somme qu'est le \italique{Dictionnaire des antiquités grecques et romaines} de Daremberg et Saglio} : bandelette que les femmes enroulaient au niveau de leur poitrine pour la soutenir -- bref : un soutien-gorge --, en revanche, le lécythe a étonnamment égaré les commentateurs: ils lui ont donné son sens (courant) de \italique{fiole à huile}. D'aucuns ont même déployé force énergie pour tenter de montrer en quoi cet accessoire était typiquement masculin (!?) 
Ils eussent plus vite et mieux fait d'ouvrir n'importe quel lexique\nc{Les dictionnaires classiques : Bailly, \LSJ, \&c. Hésychius (λ 857): τὸ μεταξὺ τοῦ λαυκανίου καὶ αὐχένος ήχῶδες. Cf. le verbe ληκυθίζειν = λαρυγγίζειν = φαρυγγίζειν, \italique{donner de la voix}.} : ils y auraient appris que λήκυθος peut aussi désigner la \italique{pomme d'Adam}. On comprend mieux l'effarement de Mnésiloque face à cette bandelette toute féminine entourant la poitrine désespérément plate d'Agathon, surmontée d'une proéminente et peu ragoûtante pomme d'Adam ; admettez que cela est d'assez mauvais goût (ὡς οὐ ξύμφορον)\nc{Faut-il en conclure que la présence d'une pomme d'Adam était caractéristique du masque comique représentant Agathon ?}.}
\comm{140}{κατόπτρου.}{La variante κάτροπτον (avec métathèse du ρ) est la seule attestée par l'épigraphie pour les inscriptions contemporaines de la représentation des \italique{Thesmophories}\nc{Voir par exemple les \italique{Inscriptiones Graecae} XXII.1471.47, 1544.58.
Le lecteur pourra aisément vérifier ces dires par lui-même en papillonnant sur les \italique{Inscriptiones Graecae} à l'adresse suivante : \href{https://epigraphy.packhum.org/}{https://epigraphy.packhum.org/}}. Ce n'est qu'à partir de la fin du \nieme{IV} siècle que la forme κάτοπτρον, plus conforme à l'étymologie, semble prendre le dessus.}
\comm{140}{τίς δαὶ κατόπτρου καὶ ξίφους κοινωνία;}{Le sel de ce vers ne se peut goûter qu'à condition d'avoir en tête le proverbe grec\nc{\nieme{3} vers d'un fragment comique adespote cité par Stobée (IV, 30, 6a) ; c'est le \nieme{903} fragment des \italique{Poetae Comici Graeci, Adespota}, VIII (1995).} :
\citComm{τίς γὰρ κατόπτρῳ καὶ τυφλῷ κοινωνία;\\
Qu'y a-t-il de commun entre un aveugle et un miroir ?}
Ce proverbe a été répertorié par les parémiographes\nc{\italique{Corpus paroemiographorum Graecorum}, II, p. 774 (\italique{Mantissae Proverbiorum} 100).} dans sa version concise (κάτοπτρον ὁ τυφλός) avec d'autres associations tout aussi inutiles. Ainsi : φαλακρὸς κτένα , \italique{un chauve un peigne}.
L'on trouvait bien évidemment des proverbes équivalents dans notre vieux parler ; \italique{À l'aveugle ne duit} [$=$ convient]\italique{ peinture / Couleur, miroir ne figure}\nc{\italique{Le livre des proverbes français} de Le Roux de Lincy, Série n°V, à l'entrée \italique{Aveugle}. À l'initiative du génial Pierre Boutang, le livre a été réédité en 1996 chez Hachette, augmenté d'une sienne préface dans laquelle -- entre autres merveilles -- Boutang questionne avec bonheur l'étymologie du \italique{pro-verbe}.}.
Aristophane joue par paronomase entre ξίφους et τυφλοῦ.
Au vers 138, il devait sans doute se livrer à un détournement du même genre : on recense de nombreux proverbes mettant en relation une lyre avec un âne (ex. : τί γὰρ κοινὸν ... λύρᾳ καὶ ὄνῳ, voir Tosi n° 1934), animal lourd et rétif aux muses ; ou une lyre avec un choucas (\latin{Nihil graculo cum fidibus}, \italique{Adages} d'Érasme n° 337), oiseau criard.}
\comm{153}{κελητίζεις, ὅταν Φαίδραν ποιῇς.}{Allusion à la \italique{pendula Venus} des Latins. Cf. \italique{Guêpes}, 500 ; \italique{Lysistrata}, 60 et 773. Le poète Machon disait καθιππάζειν (Athénée, XIII, p. 581d). (Note d'Alphonse Willems. Phèdre était éprise d'Hippolyte, dont le nom signifie littéralement le «~délieur de chevaux~» ; c'est-à-dire celui dont les chevaux courent à bride abattue.)}
\comm{157}{ὅταν σατύρους τοίνυν ποιῇς.}{Lors des Grandes Dionysies, chaque poète tragique en lice devait présenter une tétralogie composée de trois tragédies et d'un drame satyrique.}
\comm{158}{συμποιῶ σοὔπισθεν ἐστυκὼς}{; ἐστυκὼς de στύω, \italique{bander}, alors qu'on attendait ἑστηκώς, \italique{se tenir debout} ; σοὔπισθεν est mis ici pour σοί ὄπισθέν (le datif σοί réclamé par le συμ- de συμποιῶ).}
\comm{161–3}{Ἴβυκος ἐκεῖνος κἀνακρέων ὁ Τήιος / κἀλκαῖος, οἳ περὶ ἁρμονίαν ἐχύμησαν / ἐμιτροφόρουν τε καὶ διεκίνων Ἰωνικῶς,}{Les trois poètes lyriques Ibycos de Rhégium, Anacréon de Téos et Alcée de Lesbos sont aussi rapprochés par le scholiaste de Pindare (\italique{Isthmiques}, 2.1b\nc{Au tout début de la célèbre Isthmique où Pindare affirme qu'autrefois la Muse ne se faisait pas mercenaire mais chantait la beauté. Le scholiaste de noter que cela s'applique aux trois poètes lyriques : ταῦτα δὲ τείνει καὶ εἰς τοὺς περὶ Ἀλκαῖον καὶ Ἴβυκον καὶ Ἀνακρέοντα... ainsi qu'à tous les poètes apparemment fort portés sur la pédérastie : καὶ εἴ τινες τῶν πρὸ αὐτοῦ δοκοῦσι περὶ τὰ παιδικὰ ἠσχολῆσθαι· οὗτοι γὰρ παλαιότεροι Πινδάρου. Et d'appuyer le trait avec cette anecdote : on interrogeait Anacréon sur la raison qu'il avait d'écrire des hymnes à de jeunes adolescents plutôt qu'aux dieux. Anacréon répondit que c'était justement ces derniers leurs dieux : Ἀνακρέοντα γοῦν ἐρωτηθέντα, φασί, διατί οὐκ εἰς θεοὺς ἀλλ εἰς παῖδας γράφεις τοὺς ὕμνους; εἰπεῖν, ὅτι οὗτοι ἡμῶν θεοί εἰσιν.}) pour d'autres raisons que d'ordre harmonique : ce sont trois auteurs pédérastiques ; Cicéron les regroupe aussi sur le critère de la passion amoureuse (\italique{Tusculanes}, IV, 33) :
\citComm{\latin{quid denique homines doctissimi et summi poëtae de se ipsis et carminibus edant et cantibus ? Fortis vir in sua re publica cognitus quae de iuvenum amore scribit Alcaeus ! nam Anacreontis quidem tota poësis est amatoria. Maxume vero omnium flagrasse amore Reginum Ibycum apparet ex scriptis.}\\\francais{Avec quelle licence de très savants hommes et de grands poètes n'ont-ils pas chanté leurs galanteries ? Alcée, ce fameux guerrier de Mitylène, que n'a-t-il pas écrit de ces inclinations à la Grecque ? Anacréon respire-t-il autre chose que l'amour ? On voit la passion encore poussée bien plus loin dans les poésies d'Ibycus.}}Le spectateur du \nieme{V} siècle ne s'y trompait pas et savait à quoi s'en tenir avec ces trois bougres-là et comment interpréter correctement les euphémismes agathoniens.
Le commentaire ancien le plus pertinent sur notre passage n'est pas écrit en marge du \italique{Ravennas} -- comme c'est hélas bien souvent le cas pour les \italique{Thesmophories} -- mais se trouve dans la \italique{Souda} (ε 989 -- et quasiment à l'identique dans le lexique de Zonaras, p. 707) :\citComm{ἀντὶ τοῦ μίτραις\nc{Zonaras -- ou le \italique{pseudo}-Zonaras pour les esprits chagrins -- comporte un accusatif en place de l'ablatif : ἀντὶ τοῦ μίτρας ἤτοι ζώνας διαζωσάμενοι.} τουτέστι ζώναις διαζωσάμενοι ἁβρῶς περιεπάτουν.\\
Signifie : ceints de mitres -- c'est-à-dire de ceintures -- ils vivaient de façon luxurieuse.}Ce commentaire est doublement intéressant : d'abord, il précise que la μίτρα dont il est question est à prendre au sens de \italique{ceinture} et non de bandeau. Ensuite parce qu'il glose trois des mots du vers 163 : ἐμιτροφόρουν par μίτραις διαζωσάμενοι ; Ἰωνικῶς par ἁβρῶς\nc{Voir Érasme, \italique{Adages} (n° 469, \latin{Risus Ionicus}).} ; et -- plus intéressant -- le \italique{contra metrum} †διεκίνων† (†διεκίνουν† dans la \italique{Souda}) par περιεπάτουν (ici avec le sens figuré de \italique{se conduire}, \italique{vivre de telle} ou \italique{telle manière}). 
Cette glose est donc incompatible avec l'émendation de Toup, très largement reprise dans les éditions modernes. Toup proposa la correction suivante : διεκλῶντ' (de διακλάω , \italique{briser en deux}, d'où \italique{briser une vie par la mollesse} (?!)).
Il est bien préférable -- tant du point de vue paléographique que du sens -- de lire διέκειντ', de διάκειμαι, \italique{être disposé}, synonyme de περιπατέω, et qui se construit fréquemment avec un adverbe (ici : διάκειμαι Ἰωνικῶς, \italique{avoir des dispositions à la vie ionienne}, c'est-à-dire à la mollesse et la luxure).\\
Les sous-entendus d'Agathon sont à peine voilés : les trois poètes pédérastes portent la ceinture à la mode ionienne, c'est-à-dire lâche, voire franchement dénouée...
La métaphore est la suivante : tels qu'ils portaient la ceinture, béante, telles étaient leurs \italique{harmonies}. Le vers 162, οἳ περὶ ἁρμονίαν †ἐχύμησαν†\nc{La rêverie du scholiaste reprise par Suidas, ἐχύμισαν· ἀντὶ τοῦ ἔγχυμον ἐποίησαν, n'est évidemment pas à prendre au sérieux. Pas plus que l'invention du verbe χυμίζω... à moins qu'on ne m'explique à quelle sauce (χυμός) tourner le passage pour qu'il ait un rapport avec les inflexions de voix (\italique{harmonies}). Ces divagations auront probablement troublé davantage encore les scribes, déjà en butte à un texte corrompu.}, se laisse maintenant comprendre. Il faut lire :\citComm{οἳ περὶ ἁρμονίαν ἐκεχήνεσαν}La correction est très légère : une nouvelle faute de iotacisme (entre η et υ), de confusion μ/ν (quasiment identiques sur notre \italique{Ravennas}, comme sur beaucoup de manuscrits écrits en minuscules) et d'une haplographie bien explicable :  $\text{εκεχ} \xrightarrow{\text{saut du même au m. }} \text{ εχ}$.\\
χάσκειν signifie littéralement \italique{avoir la bouche grande ouverte}\nc{Bien souvent : \italique{être béat d'admiration}. Voir Taillardat §§472 et 578.}. En bâillant de la sorte aux corneilles, l'on pouvait émettre de légers sons, comparables aux coquilles grillant sur des chardons\nc{ἀνέχασκον εἷς ἕκαστος ἐμφερέστατα ὀπτωμέναις κόγχαισιν ἐπὶ τῶν ἀνθράκων. Aristophane, \italique{fr.} 67 K.-A.} ou à cette vieille lyre de Cratinos, le rival hors-d'âge d'Aristophane, chevrotant ses derniers et mauvais accords (\italique{Cavaliers}, v. 533)\nc{C'est bien ainsi qu'il faut interpréter τῶν θ' ἁρμονιῶν διαχασκουσῶν : « émettre des sons, ou accords » et non « aux joints béants » (en parlant de la lyre), comme on peut le lire.} :\citComm{τῶν θ' ἁρμονιῶν διαχασκουσῶν.}(Noter la collocation entre (δια-)χάσκω et ἁρμονία.)  C'est la même expression qu'on retrouve dans la bouche d'Agathon : περὶ ἁρμονίαν ἐκεχήνεσαν, « béer une harmonie », c'est-à-dire bêler -- ou chevroter, telle est l'image en français, point très éloignée du grec -- quelques paroles en musique\nc{ἁρμονία a le sens d'arrangement de paroles pour un morceau de musique dans Platon, \italique{Theæt.} 175e-176a.}. Le sel de la métaphore réside dans le double sens de περιχάσκειν. Le verbe est justement attesté au sens d'\italique{ouvrir une ceinture}\nc{Voir LSJ, I. b. Attesté par Héliodore le médecin (I\italique{er} ou \nieme{II} s. ap. J.-C). De longs fragments nous en ont été conservés par Oribase (\italique{Collectiones medicae} XLVIII, 58.5) : cf. les \italique{Œuvres d'Oribase}, t. 4, p. 320 (1862), de Bussemaker et Daremberg. À propos d'un patient à qui l'on applique un bandage, appelé ceinture : τῇ ζώνῃ \textbf{περιχάσκων} ζώννυται ὁ πάσχων, καὶ τότε ὁ εὐθὺς τελαμὼν ἀπὸ τῆς ὀσφύος ἄγεται ἐπὶ τὴν ἕδραν καὶ τὸν περίναιον, « On entoure le malade avec la ceinture, en l'\italique{ouvrant} largement ; puis on amène la bande droite des lombes au siège et au périnée » (trad. Daremberg).}.\\
Il est remarquable de noter que la tradition millénaire du texte d'Aristophane ait conservé la leçon ἐχύμησαν avec un η jusqu'au scribe du \italique{Ravennas} ; et ce en dépit de manuscrits anciens fautifs portant ἐχύμισαν et de commentaires afférents (ἔγχυμον ἐποίησαν) -- ayant pénétré jusqu'aux marges du \italique{Ravennas} -- suggérant la présence d'un verbe *χυμίζω. Preuve encore une fois que l'œuvre conservée d'Aristophane est, comme l'avait déjà noté Willems, l'une des plus fidèlement transmises avec l'\italique{Iliade} et l'\italique{Odyssée}\nc{Voir l'introduction à ses premières notes sur les \italique{Cavaliers}. L'une des raisons en est que le Comique, au contraire des Tragiques, ne fut que peu rejoué. Les premières éditions critiques de ses œuvres furent sans doute beaucoup plus aisées à constituer que celles d'un Euripide, dont beaucoup plus de versions différentes des pièces circulaient. Victor Coulon, l'éditeur aux Belles Lettres d'Aristophane, abonda dans le sens de Willems : à l’opposé du courant majoritaire de son temps, il s'est astreint à conserver le plus possible les leçons des manuscrits contre les conjectures des philologues. Ce « conservatisme », qu'il défend pour le cas d'Aristophane dans son \italique{Essai sur la méthode de la critique conjecturale appliquée au texte d’Aristophane} (1933) lui valut d’être encore considéré à la fin des années 90 par Sommerstein comme le meilleur éditeur d’Aristophane.}.}
\comm{164}{Φρύνιχος, τοῦτον γὰρ οὖν ἀκήκοας.}{Le poète Phrynicos, qui vivait au temps des guerres médiques portait naturellement le riche costume de son temps. Cf. les \italique{Cavaliers}, v. 1331-1332 :
\citComm{Ἀλλαντοπώλης — ὅδ᾽ ἐκεῖνος ὁρᾶν τεττιγοφόρας, ἀρχαίῳ σχήματι λαμπρός,\\
οὐ χοιρινῶν ὄζων ἀλλὰ σπονδῶν, σμύρνῃ κατάλειπτος.\\
\francais{Le charcutier — Le voilà qui apparait, la cigale à la chevelure, dans tout l'éclat du costume antique, fleurant, non les coquilles mais les libations de paix, et frotté avec de la myrrhe.}}Et Thucydide (I, 6) :\citComm{καὶ οἱ πρεσβύτεροι αὐτοῖς τῶν εὐδαιμόνων \{διὰ τὸ ἁβροδίαιτον\} οὐ πολὺς χρόνος ἐπειδὴ χιτῶνάς τε λινοῦς ἐπαύσαντο φοροῦντες καὶ χρυσῶν τεττίγων ἐνέρσει κρωβύλον ἀναδούμενοι τῶν ἐν τῇ κεφαλῇ τριχῶν.\\
Les vieillards athéniens de la classe aisée portaient encore, il n'y a pas longtemps, des tuniques (longues) de lin, et piquaient des cigales d'or dans leurs cheveux, ramassés et noués sur le sommet de la tête.}(Note de Willems.)\\
Le scholiaste explique l'incise d'Agathon (τοῦτον γὰρ οὖν ἀκήκοας) ainsi : ἐπεὶ Ἀθηναῖος ἦν.}
